\chapter{User Manual}
\label{chap:manual}

This manual explains how each person uses the system. Installation and hosting are covered in Appendix~\ref{app:hosting}.

\section{Roles at a glance}

\begin{center}
\begin{tabularx}{\linewidth}{@{}lX@{}}
\toprule
\textbf{Role} & \textbf{What the person does} \\
\midrule
Election Officer & Installs and configures the system, prepares the election, controls its stages, sets up booths, publishes results \\
Desk Officer & Registers students, checks them in, hands out slips and PINs \\
Voter & Votes once on an authorised booth \\
\bottomrule
\end{tabularx}
\end{center}

\section{Election Officer}

\subsection{First-time installation}

\begin{enumerate}
  \item Open the portal in a browser. On a new installation the installation wizard appears.
  \item Enter the institution details, logo, time zone, and the letters for the Voter Slip ID (default \code{SUE}).
  \item Create the Election Officer account and, optionally, a profile photo.
  \item Create the first election: title, voting period, positions and programmes. A template with the college's departments is pre-filled.
  \item Press \textbf{Install}, then scan the displayed code with an authenticator application and confirm. Two-step verification is mandatory for Election Officers.
\end{enumerate}

\subsection{Preparing the election}

\begin{itemize}
  \item \textbf{Programmes:} add programmes and departments or import a CSV with the columns \code{level}, \code{programme}, \code{department}, \code{terms}.
  \item \textbf{Candidates:} add each candidate with a photo (required) and a department chosen from the list.
  \item \textbf{Desk Officers:} create an account for each person who will staff a registration desk.
  \item \textbf{Voting booths:} generate a one-time booth token in the panel and enter it together with your authenticator code on the booth computer.
\end{itemize}

\subsection{Running the election}

Use the \textbf{Elections} page to move the election through its stages (Figure~\ref{fig:lifecycle} in Chapter~\ref{chap:design}). The scheduled date and time is for information only: voting starts when you press \textbf{Open voting}, stops for a break with \textbf{Pause voting} (and continues with \textbf{Resume voting}), and ends when you press \textbf{Close voting}. Positions and candidates are locked as soon as voting opens. Only one election can be active in an installation.

\subsection{Results}

\begin{enumerate}
  \item Close voting. Nothing is visible yet.
  \item Press \textbf{Publish results}.
  \item Open \textbf{Results}. Use \textbf{Present on projector} for the reveal, one position at a time (Space or right arrow to reveal and continue, left arrow to go back, F for full screen).
  \item Optionally press \textbf{Create public link} to share the results; the link can be deleted at any time.
  \item Download the CSV report if needed.
\end{enumerate}

Equal top votes are shown as a draw; the system never picks a winner by chance.

\subsection{Account and audit log}

Click your name at the top right to open your profile: change the photo, change the password (using the current password, or an authenticator code if it was forgotten), and manage two-step verification. The audit log lists every important action, grouped into categories that can be filtered.

\section{Polling and Desk Officers}

\begin{enumerate}
  \item \textbf{Station 1 (ID Verification):} Verify the student's physical institutional identity card.
  \item \textbf{Station 2 (Physical Voter Roll):} Record the student's physical signature on the printed roll.
  \item \textbf{Station 3 (Slip Distribution):} Hand the student their pre-printed voter slip containing their sequential Voter Slip ID (e.g.\ \code{SUE26-B1-1001}) and 4-digit voting PIN, cut from the official 6-per-A4 printable sheets.
  \item \textbf{Slip Management:} Monitor live booth sequential progress, void damaged or unissued slips, or download printable batch sheets by booth.
\end{enumerate}

\section{Voter (Kiosk Booth)}

\begin{enumerate}
  \item Go to the designated voting booth terminal. The screen displays the prefilled sequential Voter Slip ID (e.g.\ \code{SUE26-B1-1001}).
  \item Enter the 4-digit voting PIN using the on-screen virtual numpad with the mouse.
  \item View candidates with 1:1 square cropped photos, department, and bio. For every position, select one candidate or ``None of the Above'' (NOTA).
  \item Review all selections on the summary screen and press \textbf{Submit my vote}.
  \item The connected thermal printer silently prints the official physical \textbf{VVPAT} receipt slip.
  \item Inspect the printed paper slip, drop it into the sealed physical ballot box, and exit. The booth automatically resets for the next sequential voter.
\end{enumerate}

\notebox{A submitted ballot cannot be changed or looked up. Ballot records are stored completely decoupled from voter identity and timestamp correlation.}

\section{Troubleshooting summary}

\begin{center}
\small
\begin{tabularx}{\linewidth}{@{}>{\raggedright}p{5cm}X@{}}
\toprule
\textbf{Problem} & \textbf{Action} \\
\midrule
PIN cannot be issued & The student has already voted, or the slip ID is wrong \\
PIN stopped working & It works once, stays valid until it is used or voting closes, and is cancelled after five wrong entries; issue a new one at the desk (it replaces the old) \\
Booth says voting is not open & The election is not in the voting stage, or voting is paused \\
Booth returned to set-up & It was removed in the panel; generate a new token \\
Results are locked & Close voting, then publish the results \\
Lost authenticator & The system operator runs the recovery command described in the security guide \\
\bottomrule
\end{tabularx}
\end{center}
